\documentclass[10pt]{article}
\usepackage[margin=1.7cm]{geometry}
\usepackage{booktabs}
\usepackage{multirow}
\usepackage{tabularx}
\usepackage{array}
\usepackage{xcolor}
\usepackage{amsmath}
\usepackage{siunitx}
\usepackage{hyperref}
\hypersetup{colorlinks=true,linkcolor=blue,urlcolor=blue}
\setlength{\tabcolsep}{4pt}
\renewcommand{\arraystretch}{1.12}
\newcommand{\best}[1]{\textbf{#1}}
\newcommand{\run}{\textcolor{orange!80!black}{Running}}
\newcommand{\down}{\ensuremath{\downarrow}}
\newcommand{\up}{\ensuremath{\uparrow}}

\title{Night3R: Current Ablation Experiments}
\author{}
\date{July 16, 2026}

\begin{document}
\maketitle

\section{Final MVSEC Model Selection}

The final reported MVSEC model is \emph{not} the mixed-training checkpoint used
as the common starting point for controlled ablations.  The final MVSEC weight
is fixed to
\begin{center}
\path{outputs/night3r_mvsec_only_adapt_1k_from_best/last_event_dark3r.pth}.
\end{center}
All final MVSEC main-result, comparison, robustness, and visualization numbers
must use this checkpoint together with
\path{config/train_night3r_mvsec_only_adapt_1k_from_best.yaml}.  In particular,
the mixed-model MVSEC result (Raw AbsRel 0.4091) is only a cross-domain mixed
model diagnostic and must not be presented as the final MVSEC result.

\begin{table}[h]
\centering
\caption{Final reported MVSEC checkpoint.  These are the final-model results,
not the common reference values used inside the ablation tables.}
\label{tab:final-mvsec}
\begin{tabular}{llrrrr}
\toprule
Split & Event & Raw AbsRel \down & Raw $\delta_1$ \up & Median AbsRel \down & Median $\delta_1$ \up \\
\midrule
MVSEC-current & Real & \best{0.2073} & \best{0.8118} & \best{0.1949} & \best{0.8016} \\
MVSEC-current & Zero & 0.6798 & 0.2057 & 0.4216 & 0.3514 \\
MVSEC-current & Shuffle & 0.3165 & 0.5754 & 0.3122 & 0.5027 \\
\midrule
MVSEC-old & Real & \best{0.2222} & \best{0.8135} & \best{0.1922} & \best{0.7814} \\
MVSEC-old & Zero & 0.7585 & 0.1891 & 0.3625 & 0.4454 \\
MVSEC-old & Shuffle & 0.3060 & 0.5762 & 0.3099 & 0.4896 \\
\midrule
DSEC-current & Real & \best{0.1356} & \best{0.8373} & \best{0.0559} & \best{0.9826} \\
DSEC-current & Zero & 0.3101 & 0.2689 & 0.1071 & 0.9176 \\
DSEC-current & Shuffle & 0.1833 & 0.6305 & 0.0866 & 0.9188 \\
\bottomrule
\end{tabular}
\end{table}

The final checkpoint reduces MVSEC-current Raw AbsRel from 0.5880 for MASt3R
to 0.2073.  Real events also clearly outperform zero and shuffled events under
the same final weight.  The later VATS continuation is retained as a separate
experimental variant; it is not substituted for the selected final MVSEC
checkpoint.

On MVSEC-current, real events reduce Raw AbsRel by 69.5\% relative to zero
events and by 34.5\% relative to shuffled events.  On DSEC-current, the
corresponding reductions are 56.3\% and 26.0\%.  The strict ordering
real $<$ shuffle $<$ zero on both datasets shows that the final model benefits
from event presence and, more importantly, from correct event--RGB alignment.
The MVSEC warp PSNR changes by less than 0.1 dB across the three modes, so the
depth metrics are substantially more sensitive to event correctness than this
pairwise photometric diagnostic.

\begin{table}[h]
\centering
\caption{DL3DV event-path diagnostics for the selected final checkpoint.  The
full local DL3DV eval root has no metric depth ground truth or usable warp
supervision in this evaluator, so these are diagnostic rather than accuracy
metrics.}
\label{tab:final-dl3dv-control}
\begin{tabular}{lrrr}
\toprule
Event & Gate mean & Event density & Event occupancy \\
\midrule
Real & 0.01918 & 0.5803 & 0.6115 \\
Zero & 0.00000 & 0.0000 & 0.0000 \\
Shuffle & 0.01916 & 0.5803 & 0.6115 \\
\bottomrule
\end{tabular}
\end{table}

Real and shuffled DL3DV events have almost identical marginal statistics and
mean gates, as intended by the shuffle control.  Because no quality metric is
available on this root, these rows confirm execution and distribution
preservation but cannot establish that real alignment improves DL3DV geometry.

\section{Experimental Protocol}

Unless otherwise stated, all single-GPU ablations use the same initialization,
data, optimizer, and evaluation protocol.  The model is initialized from the
same 3k-step Night3R checkpoint.  The MASt3R RGB backbone and prediction head
are frozen, while the event encoder and the selected fusion adapter are trained
for 3000 steps.  Training uses BF16, batch size 1, balanced sampling over
MVSEC, the DSEC holdout training split, and DL3DV, and seed 0.  The event tensor
contains four temporal bins with separate positive and negative polarities,
giving eight input channels.  Images are resized to $512\times512$.

All main ablation results are evaluated with real events, no test-time low-light
augmentation, and 128 fixed adjacent pairs from MVSEC-current and DSEC-current.
The primary metric is raw metric-depth AbsRel; lower is better.  The same sample
selection and seed 17 are used for all evaluations.  Zero-event and shuffled-
event controls are additionally reported where available.

The earlier core BRG comparison used two GPUs but the same 3000-step budget.
It is therefore shown separately from the strict single-GPU table and should
not be numerically mixed with the single-GPU reference when computing small
deltas.

This separation is intentional: the final MVSEC checkpoint is used for the
paper's MVSEC main result, whereas module/loss/data ablations use one shared
pre-adaptation initialization so that the effect of the changed factor is not
confounded by target-domain MVSEC adaptation.

\section{Fusion and BRG Ablations}

\subsection{Core BRG Comparison}

\begin{table}[h]
\centering
\caption{Core fusion and BRG design ablations.  The first four rows use the
earlier 2-GPU 3k-step protocol.  The remaining completed rows use the 1-GPU
3k-step protocol.  All values are raw AbsRel.}
\label{tab:fusion}
\begin{tabular}{lccc}
\toprule
Variant & Protocol & MVSEC \down & DSEC \down \\
\midrule
Full bounded reliability gate (BRG) & 2-GPU & \best{0.4356} & 0.0844 \\
BRG without gate regularization & 2-GPU & 0.4358 & \best{0.0842} \\
Scalar gated cross-attention & 2-GPU & 0.5403 & 0.0866 \\
BRG with relaxed bounds & 2-GPU & Diverged & Diverged \\
\midrule
BRG without event statistics & 1-GPU & 0.4509 & 0.0665 \\
Legacy channel-wise tensor gate & 1-GPU & 0.7600 & 0.0728 \\
Replace RGB token by bounded event residual & 1-GPU & 0.2266 & 0.1986 \\
Concatenate RGB/event tokens + MLP residual & 1-GPU & 0.5499 & \best{0.0626} \\
Add projected event token to RGB token & 1-GPU & 0.5762 & 0.1028 \\
\bottomrule
\end{tabular}
\end{table}

The scalar cross-attention baseline computes an event-conditioned attention
update but applies one global learned scalar to every sample and token.  Its
gate standard deviation is exactly zero.  In contrast, BRG predicts a spatial
token gate from event reliability statistics.  Full BRG reduces MVSEC AbsRel
from 0.5403 to 0.4356, a relative reduction of 19.4\%, while also improving
DSEC from 0.0866 to 0.0844.  This gap is largest on MVSEC, where event sparsity
and local reliability vary strongly across the image.

Removing gate regularization changes little under the conservative 3000-step
setting because the learned gate remains below its regularization caps.  Thus,
this experiment supports describing the regularizer as a stability safeguard,
not as the main source of accuracy.  Relaxing the residual bounds causes the
event scale to grow from approximately 0.0309 to 2.57 and produces non-finite
depth.  This directly supports the bounded-residual design: the bounds protect
the pretrained MASt3R token distribution from an uncontrolled event update.

The legacy tensor baseline is not a concat baseline.  It concatenates features
only to predict a per-channel sigmoid gate and then injects gated event tokens.
The two explicit conventional baselines use the same shared initialization and
an initial residual scale of 0.035.  Concatenation followed by an MLP obtains
0.5499 on MVSEC and 0.0626 on DSEC, exposing a strong cross-domain tradeoff:
it improves DSEC but is substantially worse than BRG on MVSEC.  Unconditional
additive projection is worse on both splits (0.5762/0.1028).  These results
support reliability-conditioned bounded injection over a fixed additive event
residual; they do not support a universal DSEC advantage for BRG over concat.

\subsection{Event-Statistic and Source Calibration}

\begin{table}[h]
\centering
\caption{BRG reliability-statistic and source-calibration ablations under the
single-GPU 3k-step protocol.}
\label{tab:stats}
\begin{tabular}{lcc}
\toprule
Variant & MVSEC AbsRel \down & DSEC AbsRel \down \\
\midrule
Full BRG, seed 0 & 0.4493 & 0.0687 \\
No event-stat conditioning & 0.4509 & \best{0.0665} \\
Density only & 0.4513 & 0.0686 \\
Occupancy only & 0.4530 & 0.0682 \\
Polarity balance only & 0.4524 & 0.0680 \\
Density + occupancy, no polarity & 0.4518 & 0.0690 \\
No source bias or source-stat normalization & \best{0.4401} & 0.0718 \\
No source bias only & 0.4403 & 0.0714 \\
No source-stat normalization only & 0.4511 & 0.0687 \\
\bottomrule
\end{tabular}
\end{table}

The three event statistics are log density, occupancy, and polarity balance.
Their individual 3k-step differences are small, so the current evidence does
not justify claiming that one component independently dominates.  Source-aware
calibration shows a cross-domain tradeoff: removing source bias improves MVSEC
but degrades DSEC.  Its main purpose is therefore mixed-dataset calibration,
not a universal per-dataset gain.

\subsection{Event-Input Controls}

\begin{table}[h]
\centering
\caption{Final MVSEC-only 1k checkpoint under real, zero, and shuffled events.
Shuffling preserves the event value distribution while breaking
spatial/temporal alignment.}
\label{tab:event-control}
\begin{tabular}{llrrrr}
\toprule
Dataset & Event & Raw AbsRel \down & Raw $\delta_1$ \up & Median AbsRel \down & Median $\delta_1$ \up \\
\midrule
\multirow{3}{*}{DSEC} & Real & \best{0.1356} & \best{0.8373} & \best{0.0559} & \best{0.9826} \\
 & Zero & 0.3101 & 0.2689 & 0.1071 & 0.9176 \\
 & Shuffle & 0.1833 & 0.6305 & 0.0866 & 0.9188 \\
\midrule
\multirow{3}{*}{MVSEC} & Real & \best{0.2073} & \best{0.8118} & \best{0.1949} & \best{0.8016} \\
 & Zero & 0.6798 & 0.2057 & 0.4216 & 0.3514 \\
 & Shuffle & 0.3165 & 0.5754 & 0.3122 & 0.5027 \\
\bottomrule
\end{tabular}
\end{table}

Real events outperform both controls on both datasets.  Relative to zero
events, real events reduce raw AbsRel by 56.3\% on DSEC and 69.5\% on MVSEC;
relative to shuffled events, the reductions are 26.0\% and 34.5\%.  Real and
shuffled inputs have nearly identical event density, occupancy, and mean gate
values.  This is expected because the deterministic circular spatial/bin shift
preserves every event value while breaking event--RGB alignment.  The accuracy
gap therefore measures aligned event content rather than event count or average
fusion strength.

\paragraph{Trajectory controls.}
For MVSEC and DL3DV, relative poses are recovered by trimmed Sim(3)
registration of bidirectional point maps.  Both MASt3R head directions are
tested and the lower-residual direction is retained.  Translation is normalized
per edge, 32 adjacent-frame poses are chained, and EVO applies one Sim(3)
alignment per scene.  Values below are means of per-scene RMSE over six scenes.
DSEC is unavailable because its local test export has no pose ground truth.
These trajectory numbers use stored RGB without test-time low-light
augmentation.  MVSEC is a real night sequence; the DL3DV export records
\texttt{image\_mode=copy} and therefore uses source RGB rather than synthetic
darkening.

\begin{table*}[t]
\centering
\caption{ATE and RPE for the final MVSEC-only 1k checkpoint.  Lower is better.}
\label{tab:event-control-pose}
\begin{tabular}{llrrr}
\toprule
Dataset & Input/model & ATE (m) $\downarrow$ & RPE trans. (m) $\downarrow$ & RPE rot. ($^\circ$) $\downarrow$ \\
\midrule
\multirow{4}{*}{MVSEC} & MASt3R & \textbf{0.3705} & \textbf{0.2557} & \textbf{6.44} \\
 & Real & 0.3822 & 0.2829 & 95.77 \\
 & Zero & 0.4707 & 0.3613 & 133.19 \\
 & Shuffle & 0.4992 & 0.3086 & 111.56 \\
\midrule
\multirow{4}{*}{DL3DV} & MASt3R & \textbf{0.2570} & \textbf{0.1426} & \textbf{1.60} \\
 & Real & 0.4289 & 0.3151 & 153.71 \\
 & Zero & 0.4984 & 0.3290 & 155.16 \\
 & Shuffle & 0.3053 & 0.2516 & 139.71 \\
\midrule
DSEC & \multicolumn{4}{l}{Not available: no camera-pose ground truth.} \\
\bottomrule
\end{tabular}
\end{table*}

On MVSEC, real events reduce ATE by 18.8\% relative to zero events and by
23.4\% relative to shuffled events; translational RPE is reduced by 21.7\% and
8.4\%, respectively.  DL3DV does not show the same ordering: real improves over
zero, but shuffle has lower translation error.  The large Night3R rotation RPE
and normalized point-map residuals, compared with MASt3R, show that this final
depth-adapted checkpoint does not improve pose estimation.  The pose table is
therefore a diagnostic of event-input behavior and cross-view XYZ consistency,
not evidence of a pose state of the art.

A separate deterministic DL3DV darken test was run at exposures 0.20, 0.10,
and 0.05 with shot/read noise 0.01.  MASt3R remains stable, with ATE
0.2495/0.2368/0.2489 and rotation RPE 1.76/2.19/3.11 degrees.  Night3R with
real events obtains ATE 0.4502/0.4533/0.4124 and rotation RPE
156.12/156.54/156.79 degrees; zero and shuffled controls have similarly large
rotation error and no consistent ordering.  Thus explicit darkening does not
repair the final checkpoint's cross-view XYZ inconsistency, and this pose
diagnostic is not used as positive evidence.

\section{Dynamic Loss and Supervision Ablations}

For a sample with active loss set $\mathcal{A}$, the implemented objective is
\begin{equation}
  \mathcal{L}_{\mathrm{dyn}}
  = \frac{w_{\mathrm{target}}}
  {\max(\sum_{k\in\mathcal{A}} w_k,\epsilon)}
  \sum_{k\in\mathcal{A}} w_k\mathcal{L}_k,
\end{equation}
where source-aware multipliers modify $w_k$ for MVSEC, DSEC, and DL3DV.  A term
is excluded from both sums when its required supervision is unavailable.  This
prevents the overall gradient scale from changing simply because different
datasets expose different labels.

\begin{table}[h]
\centering
\caption{Dynamic-loss and supervision ablations.}
\label{tab:loss}
\begin{tabular}{lcc}
\toprule
Variant & MVSEC AbsRel \down & DSEC AbsRel \down \\
\midrule
Full BRG, seed 0 & \best{0.4493} & 0.0687 \\
No active-loss normalization & 0.4548 & 0.0705 \\
Unit source weights & 0.4537 & 0.0742 \\
Static plain loss & 0.4523 & 0.0713 \\
No event-depth-edge loss & 0.4553 & \best{0.0672} \\
No depth-gradient loss & 0.4514 & 0.0685 \\
No feature distillation & 0.4509 & 0.0682 \\
No structure consistency & 0.4510 & 0.0691 \\
No event-photo loss & 0.4516 & 0.0689 \\
No metric-depth loss & 0.4792 & 0.5349 \\
\bottomrule
\end{tabular}
\end{table}

The largest failure occurs without metric-depth supervision, especially on
DSEC, where AbsRel increases from 0.0687 to 0.5349.  Sparse metric depth is the
main scale anchor.  Active-loss normalization and source-aware weights provide
smaller but consistent mixed-dataset benefits.  The differences from removing
individual auxiliary terms are modest at 3000 steps and should not be
overstated.

\section{Dataset and Low-Light Ablations}

\subsection{Training-Data Composition}

\begin{table}[h]
\centering
\caption{Dataset-composition ablations.}
\label{tab:data}
\begin{tabular}{lcc}
\toprule
Training data & MVSEC AbsRel \down & DSEC AbsRel \down \\
\midrule
MVSEC + DSEC + DL3DV & 0.4493 & \best{0.0687} \\
MVSEC only & \best{0.4414} & 0.0940 \\
MVSEC + DSEC, without DL3DV & 0.4450 & 0.0752 \\
\bottomrule
\end{tabular}
\end{table}

MVSEC-only training improves the in-domain MVSEC result but substantially
degrades DSEC.  Removing DL3DV has the same, weaker pattern.  DL3DV therefore
acts mainly as a cross-domain regularizer in the current mixed training setup.

\subsection{Low-Light Generation Route}

The default physically motivated route first approximates a linear sensor
domain using inverse gamma, applies exposure scaling there, adds
signal-dependent shot noise and signal-independent read noise, clips, and then
returns to sRGB with gamma encoding:
\begin{align}
 I_{\mathrm{lin}} &= I_{\mathrm{sRGB}}^{\gamma}, \\
 I_e &= e I_{\mathrm{lin}},\qquad
 e\sim\mathrm{LogUniform}(0.03,1.0), \\
 \widetilde I_{\mathrm{lin}} &= I_e
 + \sigma_s\sqrt{I_e+\epsilon}\,\eta_s
 + \sigma_r\eta_r, \\
 I_{\mathrm{dark}} &=
 \mathrm{clip}(\widetilde I_{\mathrm{lin}},0,1)^{1/\gamma},
 \qquad \gamma=2.2.
\end{align}
This route is physically better motivated than adding Gaussian noise after
gamma compression because photon shot noise and electronic read noise originate
before the camera response curve.  It remains an approximation to RAW rather
than a claim of exact sensor simulation, since the source images have already
passed through an unknown ISP.

\begin{table}[h]
\centering
\caption{Low-light training ablations under the single-GPU 3k-step protocol.}
\label{tab:lowlight}
\begin{tabular}{p{7.7cm}cc}
\toprule
Variant & MVSEC AbsRel \down & DSEC AbsRel \down \\
\midrule
Linear domain, log-uniform exposure, shot + read noise (default) & 0.4493 & 0.0687 \\
No low-light augmentation & \best{0.2060} & 0.0940 \\
Linear domain, exposure only, no noise & 0.2081 & 0.0931 \\
Direct sRGB exposure, shot + read noise & 0.3918 & \best{0.0649} \\
Linear domain, fixed exposure $e=0.1$, shot + read noise & 0.4542 & 0.0747 \\
Linear domain, no exposure darkening, shot + read noise & 0.4323 & 0.0659 \\
Direct sRGB exposure, no noise & 0.2127 & 0.0859 \\
Linear domain, exposure + read noise only & 0.4445 & 0.0678 \\
Linear domain, exposure + shot noise only & 0.3024 & 0.0786 \\
\bottomrule
\end{tabular}
\end{table}

The completed factorial evidence shows that exposure scaling without noise
behaves similarly to no augmentation on MVSEC, while noise is important for
DSEC robustness.  Direct sRGB exposure without noise obtains 0.2127/0.0859,
close to the no-augmentation and linear exposure-only variants.  Adding shot
and read noise in sRGB improves DSEC to 0.0649 but degrades MVSEC to 0.3918.
In the linear route, read-only noise (0.4445/0.0678) nearly reproduces the
default shot-plus-read result (0.4493/0.0687), whereas shot-only noise shifts
the tradeoff toward MVSEC (0.3024) at the cost of DSEC (0.0786).  Thus read
noise is the main contributor to the default route's DSEC behavior, while shot
noise contributes more of the MVSEC gain.  The results support the physical
motivation of linear-domain noise, but direct sRGB darkening remains the better
empirical DSEC variant at this training budget.

\subsection{Distribution-Level Low-Light Check}

On 240 real MVSEC night images and 240 DL3DV source images, direct sRGB
multiplication at $e=0.03$ obtains the lowest selected histogram/statistics
distance (0.652), followed by direct sRGB $e=0.05$ (0.728).  Linear exposure at
$e=0.03$ obtains 1.063, clean source images obtain 1.919, and the current linear
shot/read simulation at $e=0.03$ obtains 2.122.  These statistics indicate that
a simple sRGB transform can match this particular MVSEC intensity histogram
more closely.  Histogram matching, however, is not equivalent to sensor
physical correctness: it conflates source-scene illumination, camera response,
black-level clipping, denoising, and tone mapping.  The paper should therefore
report both the distribution check and the downstream reconstruction ablation.

\section{Extreme-Low-Light Robustness}

\begin{table*}[t]
\centering
\caption{DSEC depth accuracy for the final MVSEC-only 1k checkpoint under
controlled synthetic low light.  RGB is inverse-gamma linearized, multiplied
by a fixed exposure, and corrupted by shot/read noise 0.01.  All methods receive
identical darkened RGB over 128 fixed adjacent pairs.}
\label{tab:final-dsec-darken-depth}
\begin{tabular}{clrrrr}
\toprule
Exposure & Input/model & Raw AbsRel $\downarrow$ & Raw $\delta_1$ $\uparrow$ & Median AbsRel $\downarrow$ & Median $\delta_1$ $\uparrow$ \\
\midrule
\multirow{5}{*}{0.20} & MASt3R & 0.3907 & 0.0602 & 0.1623 & 0.7494 \\
 & Dark3R & 0.5885 & 0.0002 & 0.1632 & 0.7585 \\
 & Real & 0.1611 & 0.7128 & \best{0.0710} & \best{0.9679} \\
 & Zero & \best{0.1182} & \best{0.8420} & 0.1096 & 0.8882 \\
 & Shuffle & 0.1868 & 0.6256 & 0.1410 & 0.7920 \\
\midrule
\multirow{5}{*}{0.10} & MASt3R & 0.5450 & 0.0082 & 0.2610 & 0.5263 \\
 & Dark3R & 0.6573 & 0.0000 & 0.1998 & 0.6591 \\
 & Real & 0.1727 & 0.6503 & \best{0.0796} & \best{0.9565} \\
 & Zero & \best{0.1393} & \best{0.7839} & 0.1325 & 0.8226 \\
 & Shuffle & 0.2105 & 0.5525 & 0.1684 & 0.7308 \\
\midrule
\multirow{5}{*}{0.05} & MASt3R & 0.8724 & 0.0008 & 0.4914 & 0.3486 \\
 & Dark3R & 0.8102 & 0.0000 & 0.2839 & 0.5164 \\
 & Real & \best{0.1859} & \best{0.5861} & \best{0.0934} & \best{0.9301} \\
 & Zero & 0.3479 & 0.2341 & 0.1939 & 0.7317 \\
 & Shuffle & 0.2555 & 0.4132 & 0.2028 & 0.6611 \\
\bottomrule
\end{tabular}
\end{table*}

The sample-keyed noise generator produces identical RGB statistics for all four
methods within each exposure.  At $e=0.05$, real events reduce Raw AbsRel by
78.7\% relative to MASt3R, 77.1\% relative to Dark3R, 46.6\% relative to zero,
and 27.2\% relative to shuffle.  Real events also reduce Median AbsRel by
67.1\% relative to Dark3R, 51.8\% relative to zero, and 53.9\% relative to
shuffle.  At $e=0.20$ and $e=0.10$, real events are best
after median alignment, while zero events have lower raw-scale error.  This
indicates a scale-calibration offset in the zero-event path rather than better
relative geometry.  The strongest-darkness row provides the clearest evidence
that correctly aligned events preserve depth when RGB visibility collapses.
These are controlled robustness tests, not an exact simulation of a particular
camera sensor.

\section{Current Status and Reporting Guidance}

All scheduled BRG, event-statistic, source-calibration, dynamic-loss,
dataset-composition, robustness, pair-distance, exposure-curve, fusion, and
low-light factorial jobs are complete.  The final-checkpoint DSEC darken
Real/Zero/Shuffle/MASt3R/Dark3R matrix is also complete at three fixed
exposures.  The
concat-MLP, additive projection,
direct-sRGB without noise, linear read-only noise, and linear shot-only noise
checkpoints all reached 3000 steps and were evaluated with real, zero, and
shuffled events on 128 fixed MVSEC-current and DSEC-current pairs.

In parallel, the selected final MVSEC-only 1k checkpoint was re-evaluated
with real, zero, and shuffled events on MVSEC-current, DSEC-current, and the
full DL3DV evaluation root.  Each setting uses 128 fixed adjacent pairs.  The
MVSEC/DSEC rows report metric-depth accuracy; DL3DV has no metric depth ground
truth in the local format and is therefore reported with warp RGB PSNR and
prediction-stability diagnostics.  This rerun explicitly uses
\path{outputs/night3r_mvsec_only_adapt_1k_from_best/last_event_dark3r.pth} for
all nine cells.  All nine runs are complete; the MVSEC/DSEC quality metrics are
reported in Table~\ref{tab:final-mvsec}, and the DL3DV diagnostics are reported
in Table~\ref{tab:final-dl3dv-control}.

Trajectory controls are also complete for the final checkpoint and MASt3R on
six 32-frame MVSEC and DL3DV scenes.  Table~\ref{tab:event-control-pose} reports
ATE and RPE.  DSEC is excluded from pose evaluation because the local test
export has no camera trajectory ground truth.

Small single-seed differences should be reported cautiously.  Existing
single-GPU Full-BRG repeats have MVSEC AbsRel $0.4407\pm0.0139$ and DSEC AbsRel
$0.0724\pm0.0027$ (population standard deviation over three seeds).  The
strongest current conclusions are the large real-vs-zero/shuffle event gaps,
the failure after removing metric-depth supervision, and divergence after
relaxing BRG bounds.  Differences of roughly 0.01 on MVSEC are not strong by
themselves under the measured run-to-run variation.

\end{document}
